% Document Structure — Report Class
% Demonstrates: chapters, sections, TOC, cross-references, footnotes, appendices.
% Compile with: pdflatex main.tex && pdflatex main.tex  (run twice for TOC/refs)

\documentclass[12pt, a4paper]{report}
\usepackage[utf8]{inputenc}
\usepackage[T1]{fontenc}
\usepackage{hyperref}      % Clickable links in TOC and references
\usepackage{amsmath}
\usepackage{amssymb}

% --- Page geometry ---
\usepackage[margin=2.5cm]{geometry}

\title{A Structured \LaTeX{} Document}
\author{LaTeX Tutorial}
\date{\today}

\begin{document}

% =============================================================================
% Front matter
% =============================================================================
\maketitle

\begin{abstract}
This document demonstrates the \texttt{report} class with chapters, sections,
a table of contents, cross-references, footnotes, and appendices. Compile
twice to resolve all references.
\end{abstract}

\tableofcontents
\listoffigures
\listoftables
\newpage

% =============================================================================
% Chapter 1
% =============================================================================
\chapter{Introduction}
\label{chap:intro}

This is the first chapter. The \texttt{report} class supports
\verb|\chapter|, which \texttt{article} does not.

\section{Motivation}
\label{sec:motivation}

LaTeX is the standard for academic writing\footnote{Especially in mathematics,
physics, and computer science.}. See Chapter~\ref{chap:methods} for our
methodology.

\section{Scope}
\label{sec:scope}

We cover basic document structuring features. For advanced mathematics, see the
\texttt{math\_advanced} example in this repository.

% =============================================================================
% Chapter 2
% =============================================================================
\chapter{Methods}
\label{chap:methods}

\section{Data Collection}
\label{sec:data}

Data was collected as described in Section~\ref{sec:motivation}.

\subsection{Instruments}
\label{sec:instruments}

We used the following instruments:
\begin{enumerate}
    \item Thermometer (accuracy $\pm 0.1^\circ$C)
    \item Scale (accuracy $\pm 0.01$\,g)
    \item Stopwatch (accuracy $\pm 0.001$\,s)
\end{enumerate}

\subsection{Procedure}
\label{sec:procedure}

The procedure followed three steps\footnote{Each step was repeated three times
for statistical significance.}:

\begin{enumerate}
    \item Calibrate instruments (see Section~\ref{sec:instruments}).
    \item Record measurements in Table~\ref{tab:results}.
    \item Analyze using Equation~\eqref{eq:mean}.
\end{enumerate}

\section{Analysis}
\label{sec:analysis}

The mean is computed as:
\begin{equation}
    \bar{x} = \frac{1}{n}\sum_{i=1}^{n} x_i
    \label{eq:mean}
\end{equation}

The standard deviation is:
\begin{equation}
    s = \sqrt{\frac{1}{n-1}\sum_{i=1}^{n}(x_i - \bar{x})^2}
    \label{eq:stddev}
\end{equation}

% =============================================================================
% Chapter 3
% =============================================================================
\chapter{Results}
\label{chap:results}

\begin{table}[ht]
\centering
\begin{tabular}{ccc}
    \hline
    \textbf{Trial} & \textbf{Value} & \textbf{Error} \\
    \hline
    1 & 9.81 & 0.02 \\
    2 & 9.79 & 0.03 \\
    3 & 9.82 & 0.01 \\
    \hline
\end{tabular}
\caption{Experimental results.}
\label{tab:results}
\end{table}

From Table~\ref{tab:results} and Equation~\eqref{eq:mean}:
$\bar{x} = 9.807$\,m/s\textsuperscript{2}, consistent with the accepted
value\footnote{The accepted value of $g$ at sea level is approximately
$9.81$\,m/s\textsuperscript{2}.}.

% =============================================================================
% Appendices
% =============================================================================
\appendix

\chapter{Derivations}
\label{app:derivations}

Full derivation of Equation~\eqref{eq:stddev} starts from the definition
of variance:

\begin{equation}
    \sigma^2 = \mathbb{E}[(X - \mu)^2]
\end{equation}

\chapter{Raw Data}
\label{app:data}

Additional raw data can be placed here, referenced from
Chapter~\ref{chap:results}.

\end{document}
